\documentclass[12pt]{article}
\usepackage[bottom]{footmisc} 
\renewcommand{\thempfootnote}{\arabic{mpfootnote}}

\usepackage[hidelinks]{hyperref}

\usepackage[utf8]{inputenc}

\usepackage[T1]{fontenc}
\usepackage[ttdefault=true]{AnonymousPro}

\usepackage{multicol}
\setlength{\columnsep}{1cm}

\usepackage{abc}
\usepackage{graphicx}
\usepackage{musicography}

\usepackage[margin=0.5in]{geometry} 
\geometry{letterpaper}


\usepackage{fancyhdr} 
\pagestyle{fancy} % options: empty , plain , fancy
\renewcommand{\headrulewidth}{0pt} % customise the layout...
\lhead{}\chead{}\rhead{}
\lfoot{}\cfoot{\thepage}\rfoot{}


\begin{document}

\thispagestyle{empty}

\begin{abc}[name=hastetothewedding,options={-O= +c -F single.fmt}]
X:1
T:Haste to the Wedding
M: 6/8
L: 1/8
K: D
|: dcB | "D"AFA A2f | "G"ede fdB | "D"AFA DFA |
 "A"GFG EFG |"D"AFA A2f | "G"ede fdB | "A"A2g faf | "D"d3 :|
|: dfg | "D"afa afa | "G"bgb bag | "D"fga agf | 
"A"gfg efg | "D"a3 f3 | "G"ede fdB | "A"A2g faf | "D"d3 :|
\end{abc}

\vskip 0.5in
\noindent\makebox[\textwidth][c]{%
  \begin{minipage}{7.0in}
    \setlength{\parindent}{1.5em}

Haste to the Wedding, a jig, has Irish, Scottish, and English origins at
least dating back to 1767 when it appears in the pantomime, \textbf{The
Elopement}, staged in London. The first verse of the song is:

    {\parindent0pt
      \raggedright
      ~\\
Come haste to the wedding, ye friends and ye neighbours, \\
The lovers their bliss can no longer delay; \\
Forget all your sorrows, your care, and your labours, \\
And let ev'ry heart beat with rapture to-day: \\
Ye votaries all, attend to my call, \\
Come revel in pleasures that never can cloy: \\
Come, see rural felicity, \\
Which love and innocence ever enjoy.\footnote{\url{https://tunearch.org/wiki/Annotation:Haste\_to\_the\_Wedding_(1)}} \\
~\\
}

The tune is found in America in the late 18th century music copybooks
(handwritten sheet music books) of Henry Livingston, a Major in the
3rd New York Regiment that participated in the failed attempt to take
Montreal from British control in 1775.\footnote{Loc. cit.} It appears
in Elias Howe’s publications, Musician’s Omnibus (Boston, 1850) and
School of Fife (Boston, 1851),\footnote{Loc. cit.} and was used on
both sides of the American Civil War as a military quickstep (brisk
music played during marching).\footnote{McCabe, Larry. \textbf{The
    Magic of Appalachian Fiddling}. Santorella Publications. 2005, p
  6.} Perhaps because of its military history, the tune became
especially popular in the Shenandoah Valley in Virginia in the early
20 th Century. \footnote{tunearch.org, Op. cit.} Haste to the Wedding
can be heard in Old Time Jams, Contradances, and Irish and Scottish
Sessions everywhere.

  \end{minipage}
  \vspace*{\fill}
}

\end{document}
